\section*{الفصل \ref{ch:b3:water-rock-interaction} --- التفاعل بين الماء والصخر}

\begin{solution}{exo:b3:water-rock-interaction:1}
الكاتيونات: $2\times48/40.08 + 2\times9.7/24.31 + 23/22.99 + 3.9/39.10 =
\qty{4.29}{meq/L}$؛ الأنيونات: $183/61.02 + 35.5/35.45 + 2\times48/96.06 =
\qty{5.00}{meq/L}$. الخطأ $100\times(4.29 - 5.00)/9.29 = \qty{-7.6}{\percent}$:
بعيدًا جدًّا عن 1 أو \qty{2}{\percent} في التحليل الدقيق لمثل هذا الماء.
إمّا أن يغيب أحد الكاتيونات أو يُقدَّر بأقلّ من قيمته، أو أن يكون أحد الأنيونات مرتفعًا أكثر من اللازم.
\end{solution}

\begin{solution}{exo:b3:water-rock-interaction:2}
$\sqrt I = 0.0707$. \ce{Ca^2+}: $-\log\gamma = 0.5085\times4\times0.0707/(1 +
0.3281\times6\times0.0707) = 0.126$, $\gamma = 0.748$; \ce{HCO3-}: $\gamma =
0.927$; \ce{SO4^2-}: $\gamma = 0.739$.
\end{solution}

\begin{solution}{exo:b3:water-rock-interaction:3}
$\sqrt I = 0.126$; $\gamma_{\ce{Ca}} = 0.622$, $\gamma_{\ce{SO4}} = 0.602$;
$\log\mathrm{IAP} = \log(3.0\times10^{-3}\times0.622\times4.0\times10^{-3}
\times0.602) = -5.35$؛ $\mathrm{SI} = -5.35 + 4.55 = -0.80$. الماء
غير مشبَّع ويمكنه أن يذيب الجبس.
\end{solution}

\begin{solution}{exo:b3:water-rock-interaction:4}
$m = 0.012/60.08 = \qty{2.00e-4}{mol/L}$، $\log m = -3.70$. الكوارتز:
$\mathrm{SI} = -3.70 + 3.98 = +0.28$، فوق التشبّع بقليل؛ السيليكا
اللابلورية: $-3.70 + 2.71 = -0.99$، غير مشبّعة.
\end{solution}

\begin{solution}{exo:b3:water-rock-interaction:5}
عند \qty{25}{\celsius}، $\log K = -8.45 - 6.35 - 1.47 + 10.33 = -5.94$، $m =
(10^{-5.94}\times10^{-2}/4)^{1/3} = \qty{1.42}{mmol/L}$، \qty{57}{mg/L}؛ عند
\qty{10}{\celsius}، $\log K = -5.59$، $m = \qty{1.86}{mmol/L}$،
\qty{75}{mg/L}: أكثر بنحو الثلث في البرودة.
\end{solution}

\begin{solution}{exo:b3:water-rock-interaction:6}
$[\ce{CO2}] = 10^{-7.60}\times4.0\times10^{-3}/10^{-6.35} =
\qty{2.25e-4}{mol/L}$; $[\ce{CO3^2-}] = 10^{-10.33}\times4.0\times10^{-3}/
10^{-7.60} = \qty{7.4e-6}{mol/L}$; $\log P_{\ce{CO2}} = -3.648 + 1.47 =
-2.18$، أي 16 مرة قدر الهواء الجوي. المجموع \qty{4.23}{mmol/L}: \ce{CO2}
\qty{5.3}{\percent}، \ce{HCO3-} \qty{94.5}{\percent}، \ce{CO3^2-}
\qty{0.18}{\percent}.
\end{solution}

\begin{solution}{exo:b3:water-rock-interaction:7}
$r = 10^{-12.09} = \qty{8.2e-13}{mol.m^{-2}.s^{-1}}$ (الحدود الحمضية والمحايدة
والقاعدية تُسهم بالقيم $10^{-12.45}$، $10^{-12.56}$ و$10^{-12.74}$). الشعاع
يتناقص بمعدّل $rV_m = 8.2\times10^{-13}\times1.0131\times10^{-4} =
\qty{8.3e-17}{m/s}$، \qty{2.6}{nm} في السنة؛ أمّا الشعاع \qty{0.5}{mm} فيدوم
$6.0\times10^{12}$~s، نحو \num{190000} سنة.
\end{solution}

\begin{solution}{exo:b3:water-rock-interaction:8}
$\exp[(69.8/0.008314)(1/278.15 - 1/298.15)] = 7.6$؛ مع
\qty{74}{kJ/mol}، 8.6.
\end{solution}

\begin{solution}{exo:b3:water-rock-interaction:9}
بمساواة $(12.703 - X - 3Y)/3$ و$(7.435 - 2Y)/2$ يُحذف $Y$: $X =
12.703 - 11.153 = 1.55$، أي خط أفقي. لدى الماء $X = \log(2.0\times
10^{-3}/39.10) + 7.0 = 2.71$ و$Y = \log(0.015/60.08) = -3.60$: أعلى من 1.55
وأدنى بكثير من خط الموسكوفيت--الفلسبار K ($X = -3.161 - 3Y = 7.65$)، لذا
يتكوّن الموسكوفيت.
\end{solution}

\begin{solution}{exo:b3:water-rock-interaction:10}
$\delta_w^f = -2 - 2 = -4$. مغلق: $W/R = 0.5\times(-2 - 6)/(-13 + 4) =
0.44$. مفتوح: $0.5\ln[(-8)/(-9) + 1] = 0.5\ln1.889 = 0.32$. في النظام المفتوح تغادر كلّ
كتلة مائية في حالة توازن مع الصخر الذي لاقته أخيرًا، حاملةً أكبر قدر ممكن من
الأكسجين-18؛ أمّا في النظام المغلق فتنتهي الدفعة كلّها عند $-4$
وتأخذ منه أقلّ لكلّ كيلوغرام.
\end{solution}

\begin{solution}{exo:b3:water-rock-interaction:11}
\omterm{def:b3:water-rock-interaction:open-closed}{نظام الكربونات المفتوح}: $\log K = -5.59$، $m = (10^{-5.59}\times10^{-1.5}/4)^{1/3} =
\qty{2.73}{mmol/L}$، \qty{109}{mg/L}. \omterm{def:b3:water-rock-interaction:open-closed}{نظام الكربونات المغلق}: $C_0 = 10^{-1.27 - 1.5} =
\qty{1.70e-3}{mol/L}$ و$K_{\mathrm{cal}}K_1/K_2 = 10^{-4.32}$؛ المعادلة التكعيبية
$4m^3 = 4.79\times10^{-5}(1.70\times10^{-3} - m)$ تعطي $m =
\qty{1.45}{mmol/L}$، \qty{58}{mg/L}، نحو النصف. ثاني أكسيد الكربون المتبقّي هو
$0.25\times10^{-3}$، \qty{15}{\percent} من $C_0$، لذا ينخفض $\log P_{\ce{CO2}}$
من $-1.5$ إلى $-2.33$.
\end{solution}

\begin{solution}{exo:b3:water-rock-interaction:12}
$1 - 10^{\mathrm{SI}}$: 0.99، 0.90 و0.21. قانون القوّة يعطي
$(10^4/6)^{-0.61} = 0.011$: معدّل أقلّ بنحو 90 مرة. المعدن الحقلي
قديم، زالت أسطحه التفاعلية، ومغمور بماء مسامي بطيء قريب من
التوازن؛ يتضاعف أثر العاملين معًا، ويهبط المعدّل الحقلي بعدّة
مراتب من القدر دون المعدّل المخبري.
\end{solution}

\begin{solution}{pb:b3:water-rock-interaction:1}
\textbf{1.} \ce{Ca} 2.795, \ce{Mg} 0.165, \ce{Na} 0.305, \ce{K} 0.031,
\ce{HCO3} 5.572, \ce{Cl} 0.310, \ce{SO4} 0.146, \ce{NO3} \qty{0.145}{mmol/L}.

\textbf{2.} الكاتيونات \qty{6.25}{meq/L}، الأنيونات \qty{6.32}{meq/L}: الخطأ
\qty{-0.5}{\percent}، مقبول.

\textbf{3.} $I = \tfrac12(4\times2.795 + 4\times0.165 + 0.305 + 0.031 + 5.572
+ 0.310 + 4\times0.146 + 0.145)\times10^{-3} = 0.0094$.

\textbf{4.} $\sqrt I = 0.0969$: $\gamma_{\ce{Ca^2+}} = 0.683$,
$\gamma_{\ce{HCO3-}} = 0.904$, $\gamma_{\ce{CO3^2-}} = 0.676$.

\textbf{5.} نعم: $I$ يساوي عُشر الحدّ الذي يقارب 0.1.

\textbf{6.} $a_{\ce{HCO3-}} = 5.572\times10^{-3}\times0.904 = 5.04\times
10^{-3}$ ($\log = -2.298$).

\textbf{7.} $a_{\ce{CO2}} = 10^{-7.20}\times5.04\times10^{-3}/10^{-6.46} =
9.17\times10^{-4}$ ($\log = -3.038$).

\textbf{8.} $\log P_{\ce{CO2}} = -3.038 + 1.27 = -1.77$, \qty{0.017}{atm}:
$10^{1.60} = 40$ مرّة قدر الهواء، ضمن مدى المياه الجوفية. ويأتي ثاني
أكسيد الكربون من هواء التربة في منطقة التغذية، ومن الجذور والتحلّل.

\textbf{9.} $a_{\ce{CO3^2-}} = 10^{-10.49}\times5.04\times10^{-3}/10^{-7.20}
= 2.58\times10^{-6}$، $[\ce{CO3^2-}] = \qty{3.8e-6}{mol/L}$. المجموع
\qty{6.49}{mmol/L}: \ce{CO2} \qty{14.1}{\percent}، \ce{HCO3-}
\qty{85.8}{\percent}، \ce{CO3^2-} \qty{0.06}{\percent}.

\textbf{10.} القلوية $5.572 + 2\times0.0038 = \qty{5.580}{meq/L}$:
بنسبة \qty{99.9}{\percent} من البيكربونات.

\textbf{11.} $a_{\ce{Ca^2+}} = 2.795\times10^{-3}\times0.683 = 1.909\times
10^{-3}$; $\log\mathrm{IAP} = \log(1.909\times10^{-3}\times2.58\times10^{-6})
= -8.307$; $\mathrm{SI} = -8.307 + 8.35 = +0.04$.

\textbf{12.} ضمن $\pm0.1$، وهو عدم اليقين في المؤشّر: الماء في
توازن مع الكالسيت. وقد مكث طويلًا في الصخر، فهو الينبوع البطيء
المنتشر.

\textbf{13.} $m = [10^{-5.59}\times0.0171/(4\times0.683\times0.904^2)]^{1/3}
= \qty{2.70}{mmol/L}$، \qty{108}{mg/L}، مقابل \qty{2.79}{mmol/L} المقيسة:
\qty{3}{\percent} أكثر، وهذا ما يعنيه مؤشّر قدره $+0.04$.

\textbf{14.} $a_{\ce{CO2}} = 10^{-1.27 - 3.37} = 2.29\times10^{-5}$;
$a_{\ce{H+}} = 10^{-6.46}\times2.29\times10^{-5}/5.04\times10^{-3}$، \babelsublr{pH 8.80}؛
$a_{\ce{CO3^2-}} = 1.03\times10^{-4}$; $\mathrm{SI} = \log(1.909\times
10^{-3}\times1.03\times10^{-4}) + 8.35 = +1.65$ (1.61 مع ميزان القلوية
الكامل للشكل).

\textbf{15.} بتكرار الحساب $m = [10^{-5.59 - 3.37}/(4\gamma_{\ce{Ca}}
\gamma_{\ce{HCO3}}^2)]^{1/3}$ مع $I = 3m$: $m = \qty{0.72}{mmol/L}$
(\qty{29}{mg/L})، $I = 0.0022$. ترسّب الكالسيت: $2.795 - 0.718 =
\qty{2.08}{mmol/L}$، أي نحو \qty{210}{mg} لكلّ لتر.

\textbf{16.} $0.208\times2.0 = \qty{0.42}{g/s}$، \qty{13}{t} في السنة،
$13.1/2.71 = \qty{4.8}{m^3}$؛ على مساحة \qty{400}{m^2}، \qty{12}{mm} في السنة،
\qty{0.03}{mm} في اليوم: أدنى بكثير من الترسّبات الكلسية للينابيع الباردة المذكورة في الفصل.

\textbf{17.} $X = \log(1.2\times10^{-3}/39.10) + 5.8 = 1.29$; $Y =
\log(9.0\times10^{-3}/60.08) = -3.82$.

\textbf{18.} $X < 1.55$ و$Y > -4.39$: الكاولينيت. الكوارتز: $\mathrm{SI} =
-3.82 + 3.98 = +0.16$.

\textbf{19.} دون $\log a_{\ce{H4SiO4}} = -4.39$، \qty{2.4}{mg/L} من
السيليكا.

\textbf{20.} $X = 1.55$ يتطلّب \mbox{\textdir TLT pH $= 1.55 + 4.51 = 6.06$}.

\textbf{21.} الحمضي $-10.06 - 0.500\times5.8 = -12.96$؛ المتعادل $-12.41$؛ القاعدي
$-21.20 + 0.823\times5.8 = -16.43$. المجموع: $\log r = -12.30$، الحدّ المتعادل
هو المهيمن.

\textbf{22.} \omterm{def:b3:water-rock-interaction:regimes}{النظام المحدود بالإمداد}: التربة السميكة والمنحدرات اللطيفة تعني تعرية بطيئة؛
الفلسبار الطازج الذي يصل إلى جبهة التجوية يُغسَل كلّيًا،
ويتبع الدفق معدّل التعرية لا المناخ.

\textbf{23.} $\exp[(74/0.008314)(1/283.15 - 1/298.15)] = 4.9$.

\textbf{24.} مؤشّر التشبّع للكالسيت لمياه الينبوع عند مخرجه
هو $\mathrm{SI} = \textbf{+0.04}$.
\end{solution}
